
\newpage
%%%%%%%%%%%%%%%%%%%%%%%%%%%%%%%%%%%%%%%%%%%%%%%%%%%%%%%%%%%%%%%%%%%%%%%%%%%%%%%%
\section{Interfaces to MATLAB, Octave, Python, Julia, Go, Java, ...} %%%%%%%%%%%
%%%%%%%%%%%%%%%%%%%%%%%%%%%%%%%%%%%%%%%%%%%%%%%%%%%%%%%%%%%%%%%%%%%%%%%%%%%%%%%%

The MATLAB/Octave interface to SuiteSparse:GraphBLAS is included with this
distribution, described in Section~\ref{octave}.
Python, Julia, Go, and Java interfaces are available.
These are not part of the SuiteSparse:GraphBLAS distribution.
See the links below.

%===============================================================================
\subsection{MATLAB/Octave Interface}
%===============================================================================
\label{octave}

An easy-to-use MATLAB/Octave interface for SuiteSparse:GraphBLAS is available;
see the documentation in the \verb'GraphBLAS/GraphBLAS' folder for details.
Start with the \verb'README.md' file in that directory.  An easy-to-read output
of the MATLAB demos can be found in \verb'GraphBLAS/GraphBLAS/demo/html'.

The quick start instructions on Linux or Windows are to type the following in
MATLAB/Octave (assuming you have GraphBLAS installed in your home directory).
See the \verb'README.md' for additional required instructions if you have a Mac.

{\footnotesize
    \begin{verbatim}
        cd ~/GraphBLAS/GraphBLAS
        graphblas_install
        addpath (pwd)
        cd test
        gbtest \end{verbatim}}

The MATLAB/Octave interface provides two classes that hold a matrix, vector or
scalar: \verb'GrB' value class, and the \verb'GhB' handle class.  The latter
is new to GraphBLAS 10.4.0.  It exploits the memory arenas and is much faster
than using \verb'GrB' value class.  The two classes can be mixed arbitrarily,
along with MATLAB/Octave built-in matrices.  For details, type the following
in the MATLAB/Octave command window:

{\footnotesize
    \begin{verbatim}
        help GrB
        help GhB \end{verbatim}}

Both objects contain a GraphBLAS matrix, either double or single precision
(real or complex), boolean, or any of the built-in integer types.
MATLAB/Octave sparse and full matrices can be arbitrarily mixed with GraphBLAS
matrices.  The following overloaded operators and methods all work as you would
expect for any matrix.  The matrix multiplication \verb'A*B' uses the
conventional \verb'PLUS_TIMES' semiring.

{\footnotesize
\begin{verbatim}
    A+B    A-B   A*B    A.*B   A./B   A.\B   A.^b    A/b    C=A(I,J)
    -A     +A    ~A     A'     A.'    A&B    A|B     b\A    C(I,J)=A
    A~=B   A>B   A==B   A<=B   A>=B   A<B    [A,B]   [A;B]  A(1:end,1:end) \end{verbatim}}

The new \verb'GhB' handle object is faster than the \verb'GrB' value object
because the former can be modified in-place with calls to \verb'GhB' methods,
and it can exploit the non-blocking feature of GraphBLAS.  The non-blocking
mode of GraphBLAS allows work to be left pending in a matrix, to be done later
when it is more efficient.  This cannot be done with the \verb'GrB' value
object.

All methods work on any mix of \verb'GrB', \verb'GhB', and/or MATLAB/Octave
built-in matrices.  A few new methods work only on \verb'GhB' methods; see the
MATLAB \verb'help GhB' for details.

The only minor downside to the \verb'GhB' handle object is that it changes the
MATLAB/Octave evaluation of a simple assignment statement, \verb'C=A'.  If
\verb'A' is a \verb'GrB' value object, \verb'C' is created as a copy of
\verb'A', and changing either one later on does not change the other.  If
instead \verb'A' is a \verb'GhB' handle object, \verb'C=A' does not make a
copy.  Instead \verb'C' is now just another name of the \verb'A' matrix.
Changing \verb'A' causes \verb'C' to change, and visa versa.  To construct
\verb'C' as a deep copy of the \verb'GhB' object \verb'A', use \verb'C=GhB(A)'
instead of \verb'C=A'.  This is a built-in feature of any handle class in
MATLAB/Octave that cannot be changed.  The handle object characteristics are
required for the \verb'GhB' matrix to hold pending work, but this comes with a
change in the semantics of \verb'C=A' that cannot be avoided.

Internally, the two classes require the use of different memory arenas,
introduced in GraphBLAS 10.4.0 and described in Section~\ref{sec:arenas}.  The
\verb'GrB' object uses \verb'mxMalloc' and \verb'mxFree', while the \verb'GhB'
handle object uses the system \verb'malloc' and \verb'free'.  The use of
different arenas are not visible to user application code written in
MATLAB/Octave, however.

{\bf Limitations:}
Some features for MATLAB/Octave matrices are not yet available for
GraphBLAS matrices.  Some of these may be added in future releases.

\begin{packed_itemize}
    \item The MATLAB \verb'whos' command:  \verb'GrB' matrices with dimension
        larger than \verb'2^53' do not display properly in the \verb'whos'
        command.  The dimensions are displayed correctly with \verb'disp' or
        \verb'display'.  The size in bytes of \verb'GhB' matrices are not
        displayed properly by the \verb'whos' command; they just show 8 bytes
        regardless of the size of the matrix.  Use \verb'GhB.bytes(A)' or
        \verb'GrB.bytes(A)' to obtain the true size in bytes of any matrix A
        (\verb'A' can be a built-in MATLAB/Octave matrices, \verb'GrB'
        matrices, or \verb'GhB' matrices).
    \item Linear indexing: \verb'A(:)' for a 2D matrix, and \verb'I=find(A)'.
    \item Singleton expansion.
    \item Dynamically growing arrays, where \verb'C(i)=x' can increase
        the size of \verb'C'.
    \item Saturating element-wise binary and unary operators for integers.
        For \verb'C=A+B' with MATLAB \verb'uint8' matrices, results
        saturate if they exceed 255.  This is not compatible with
        a monoid for \verb'C=A*B', and thus MATLAB does not support
        matrix-matrix multiplication with \verb'uint8' matrices.
        In GraphBLAS, \verb'uint8' addition acts in a modulo fashion.
    \item Solvers, so that \verb'x=A\b' could return a GF(2) solution,
        for example.
    \item Sparse matrices with dimension higher than 2.
\end{packed_itemize}


%===============================================================================
\subsection{Performance of MATLAB versus GraphBLAS}
%===============================================================================
\label{matlab_performance}

MATLAB R2021a uses SuiteSparse:GraphBLAS as a built-in library, but
uses it only for \verb'C=A*B' when both \verb'A' and \verb'B' are sparse.  In
prior versions of MATLAB, \verb'C=A*B' relied on the \verb'SFMULT' and
\verb'SSMULT' packages in SuiteSparse, which are single-threaded (also written
by this author).  The GraphBLAS \verb'GrB_mxm' is up to 30x faster on a 20-core
Intel Xeon, compared with \verb'C=A*B' in MATLAB R2020b and earlier.  With
MATLAB R2021a and later, the performance of \verb'C=A*B' when using MATLAB
sparse matrices is identical to the performance for GraphBLAS matrices, since
the same code is being used by both (\verb'GrB_mxm').

Other methods in GraphBLAS are also faster, some {\em extremely} so, but are
not yet exploited as built-in operations MATLAB.  In particular, the statement
\verb'C(M)=A' (where \verb'M' is a logical matrix) takes under a second for a
large sparse problem when using GraphBLAS via its \verb'GrB' interface.  By
stark contrast, MATLAB would take about 4 or 5 days, a speedup of about
500,000x.  For a smaller problem, GraphBLAS takes 0.4 seconds while MATLAB
takes 28 hours (a speedup of about 250,000x).  Both cases use the same
statement with the same syntax (\verb'C(M)=A') and compute exactly the same
result.  Below are the results for \verb'n'-by-\verb'n' matrices in GraphBLAS
v5.0.6 and MATLAB R2020a, on a Dell XPS13 laptop (16GB RAM, Intel(R) Core(TM)
i7-8565U CPU @ 1.80GHz with 4 hardware cores).  GraphBLAS is using 4 threads.

\vspace{0.10in}
{\scriptsize
\begin{tabular}{rrr|rrr}
\hline
\verb'n'    & \verb'nnz(C)' & \verb'nnz(M)' & GraphBLAS (sec) & MATLAB (sec) & speedup \\
\hline
2,048        & 20,432         & 2,048          & 0.005     & 0.024     & 4.7 \\
4,096        & 40,908         & 4,096          & 0.003     & 0.115     & 39 \\
8,192        & 81,876         & 8,191          & 0.009     & 0.594     & 68 \\
16,384       & 163,789        & 16,384         & 0.009     & 2.53      & 273 \\
32,768       & 327,633        & 32,767         & 0.014     & 12.4      & 864 \\
65,536       & 655,309        & 65,536         & 0.025     & 65.9      & 2,617 \\
131,072      & 1,310,677      & 131,070        & 0.055     & 276.2     & 4,986 \\
262,144      & 2,621,396      & 262,142        & 0.071     & 1,077     & 15,172 \\
524,288      & 5,242,830      & 524,288        & 0.114     & 5,855     & 51,274 \\
1,048,576    & 10,485,713     & 1,048,576      & 0.197     & 27,196    & 137,776 \\
2,097,152    & 20,971,475     & 2,097,152      & 0.406     & 100,799   & 248,200 \\
4,194,304    & 41,942,995     & 4,194,304      & 0.855  & 4 to 5 days? & 500,000?\\
\hline
\end{tabular}}
\vspace{0.10in}

The assignment \verb'C(I,J)=A' in MATLAB, when using \verb'GrB' objects, is up
to 1000x faster than the same statement with the same syntax, when using MATLAB
sparse matrices instead.  Matrix concatenation \verb'C = [A B]' is about 17
times faster in GraphBLAS, on a 20-core Intel Xeon.  For more details, see the
\verb'GraphBLAS/GraphBLAS/demo' folder and its contents.

Below is a comparison of other methods in SuiteSparse:GraphBLAS, compared with
MATLAB 2021a.  SuiteSparse:GraphBLAS: v6.1.4 (Jan 12, 2022), was used, compiled
with gcc 11.2.0.  The system is an Intel(R) Xeon(R) CPU E5-2698 v4 @ 2.20GHz
(20 hardware cores, 40 threads), Ubuntu 20.04, 256GB RAM.  Full details appear
in the \verb'GraphBLAS/GraphBLAS/demo/benchmark' folder.  For this matrix,
SuiteSparse:GraphBLAS is anywhere from 3x to 17x faster than the built-in
methods in MATLAB.  This matrix is not special, but is typical of the relative
performance of many large matrices.  Note that two of these (\verb'C=L*S' and
\verb'C=S*R') rely on an older version of SuiteSparse:GraphBLAS (v3.3.3) built
into MATLAB R2021a.

{\footnotesize
\begin{verbatim}
    Legend:
    S: large input sparse matrix (n-by-n), the GAP-twitter matrix
    x: dense vector (1-by-n or n-by-1)
    F: dense matrix (4-by-n or n-by-4)
    L: 8-by-n sparse matrix, about 1000 entries
    R: n-by-8 sparse matrix, about 1000 entries
    B: n-by-n sparse matrix, about nnz(S)/10 entries
    p,q: random permutation vectors

    GAP/GAP-twitter: n: 61.5784 million nnz: 1468.36 million
    (run time in seconds):
    y=S*x:   MATLAB:  22.8012 GrB:   2.4018 speedup:     9.49
    y=x*S:   MATLAB:  16.1618 GrB:   1.1610 speedup:    13.92
    C=S*F:   MATLAB:  30.6121 GrB:   9.7052 speedup:     3.15
    C=F*S:   MATLAB:  26.4044 GrB:   1.5245 speedup:    17.32
    C=L*S:   MATLAB:  19.1228 GrB:   2.4301 speedup:     7.87
    C=S*R:   MATLAB:   0.0087 GrB:   0.0020 speedup:     4.40
    C=S'     MATLAB: 224.7268 GrB:  22.6855 speedup:     9.91
    C=S+S:   MATLAB:  14.3368 GrB:   1.5539 speedup:     9.23
    C=S+B:   MATLAB:  15.5600 GrB:   1.5098 speedup:    10.31
    C=S(p,q) MATLAB:  95.6219 GrB:  15.9468 speedup:     6.00    \end{verbatim}
}




%===============================================================================
\subsection{Python Interface}
%===============================================================================
\label{python}

See Michel Pelletier's Python interface at
\url{https://github.com/michelp/pygraphblas};
it also appears at
\url{https://anaconda.org/conda-forge/pygraphblas}.

See Jim Kitchen and Erik Welch's (both from Anaconda, Inc.) Python interface at
\url{https://github.com/python-graphblas/python-graphblas} (formerly known as grblas).
See also \url{https://anaconda.org/conda-forge/graphblas}.

%===============================================================================
\subsection{Julia Interface}
%===============================================================================
\label{julia}

The Julia interface is at
\url{https://github.com/JuliaSparse/SuiteSparseGraphBLAS.jl}, developed by Will
Kimmerer, Abhinav Mehndiratta, Miha Zgubic, and Viral Shah.
Unlike the MATLAB/Octave interface (and like the Python interfaces) the Julia
interface can keep pending work (zombies, pending tuples, jumbled state) in
a \verb'GrB_Matrix'. This makes Python and Julia the best high-level interfaces
for SuiteSparse:GraphBLAS.  MATLAB is not as well suited, since it does not
allow inputs to a function or mexFunction to be modified, so any pending
work must be finished before a matrix can be used as input.

%===============================================================================
\subsection{Go Interface}
%===============================================================================
\label{go}

Pascal Costanza (Intel) has a Go interface to GraphBLAS and LAGraph:
\begin{itemize}
\item forGraphBLASGo: \url{https://github.com/intel/forGraphBLASGo}, which
is almost a complete wrapper for SuiteSparse:GraphBLAS.
Documentation is at \url{https://pkg.go.dev/github.com/intel/forGraphBLASGo}.
\item forLAGraphGo: \url{https://github.com/intel/forLAGraphGo}, which is in
progress.  Documentation is at
\url{https://pkg.go.dev/github.com/intel/forLAGraphGo}.
\end{itemize}

%===============================================================================
\subsection{Java Interface}
%===============================================================================
\label{java}

Fabian Murariu is working on a Java interface.
See \newline
\url{https://github.com/fabianmurariu/graphblas-java-native}.

